\scanpage{22}
\mspage{20}
tion de compatibilité avec le changement de base, on constate que $\operatorname{DR}(f_p)^{(p)}$ \deleted{\illegible{1}} n'est autre que $\operatorname{DR}(f_p^{(p)})$, où $f_p^{(p)}:X_p^{(p/S_p)}\longrightarrow S_p$ est le morphisme structural du frobeniusé relatif de $X_p$ sur $S_p$. Or on a le morphisme de frobenius $X_p^{(p/S)}\longrightarrow X$, qui est un $S$\uline{-}morphisme, et par fonctorialité de $\operatorname{DR}$ nous donne $\operatorname{DR}(f)_p^{(p)}\longrightarrow\operatorname{DR}(f)_p$ comme on voulait. Il faut prouver que cet homomorphisme est en fait une isogénie, donc que l\deleted{e}\inserted{'}isocristal défini par $\operatorname{DR}(f)$ devient, à l'aide des $F_p$, un bi-isocristal. Mais utilisant la relation de dualité écrite dans 2.2. (à vrai dire, l'écriture $T^d$ pour le cristal unité ne prend son sens que lorsque on le regarde comme muni de sa structure de \uuline{bi}-cristal naturelle, qui n'intervient qu'ici), on peut transposer $F$ en un $V$, tel que $FV=VF=p^{2d}$, ce qui prouve notre assertion.

D'ailleurs, lorsque l'on passe du cristal de De Rham $\operatorname{DR}(f)$ à l'isocristal correspondant, \deleted{et} donc des objets de cohomologie $\operatorname{DR}^i(f)$ aux isocristaux correspondants, \deleted{on voit que} il \deleted{doit être} \inserted{sera} vrai (tout comme pour les $\underline{R}^if_*$ de De Rham en caractéristique nulle) que leur formation commute à tout changement de base, de sorte que chacun des isocristaux $\operatorname{DR}^i(f)$ devient à son propre titre un bi-isocristal. Au moment de rédiger 1.10. il m'avait semblé \inserted{que}, sans mettre du iso dans le coup, \deleted{que} $\operatorname{DR}^i(f)$ devrait être \inserted{un} bi-cristal de poids $i$, mais je m'étais canulé, il faudrait pour cela une polarisation de $X$ sur $S$ qui définisse un isomorphisme (pas seulement une isogénie) de $\operatorname{DR}^i(f)$ avec $\operatorname{DR}^{2d-i}\otimes T^{-(d-i)}$, ce qui n'existe évidemment que très exceptionnellement; une fois qu'on l'a, on définit $V$ dans $\operatorname{DR}^i$ en transposant $F$ dans $\operatorname{DR}^{2d-i}$.

2.4. \uline{Cas d'un schéma abélien}. Il semble que ce ne soit guère que dans ce cas-là où il soit bien raisonnable de parler de \inserted{bi-}cristaux et non seulement de bi-isocristaux. De façon précise, $\operatorname{DR}^1(f)$ a dans ce cas une structure de bicristal, défini par les $F_p$ précédents, et des $V_p$ qui se définissent encore, par fonctorialité de $\operatorname{DR}$, à l'aide de l'homomorphisme ``Verschiebung'' $A_p^{(p/S_p)}\longrightarrow A_p$ (défini avec la généralité qui convient dans les notes de Gabriel dans SGAD). Comme
\margin{21}
